\documentclass[11pt]{article}
\usepackage[margin=1in]{geometry}
\usepackage[T1]{fontenc}
\usepackage{lmodern,amsmath,amssymb,amsthm,microtype}
\usepackage[colorlinks=true,linkcolor=blue,citecolor=blue,urlcolor=blue]{hyperref}
\hypersetup{pdftitle={Depth-uniform Gaussian Bayesian learning of shallow Boolean targets},pdfauthor={ChatGPT-6 Astra; Kaarel Hanni; Dmitry Vaintrob}}
\newtheorem{theorem}{Theorem}[section]
\newtheorem{proposition}[theorem]{Proposition}
\newtheorem{lemma}[theorem]{Lemma}
\newtheorem{corollary}[theorem]{Corollary}
\newcommand{\op}{\mathrm{op}}
\newcommand{\poly}{\operatorname{poly}}
\title{Depth-uniform Gaussian Bayesian learning\\of shallow Boolean targets}
\author{ChatGPT-6 Astra \and Kaarel H\"anni \and Dmitry Vaintrob\thanks{Authors are listed in alphabetical order.}}
\date{17 October 2026}
\begin{document}
\maketitle
\begin{abstract}
Consider fully connected tanh networks with an independent Gaussian prior,
conditioned on absolute output margin greater than $1/2$ on the Boolean
cube and on noiseless training labels. Suppose the target has a
one-hidden-layer representation of width $C$, weights bounded by $R$, and
signed output margin at least one. For each fixed $\zeta>0$, we prove a
sample bound polynomial in input dimension, $C,R$, learner depth, inverse
error, and logarithmic inverse confidence, simultaneously for every learner
width at least the target's actual bounded-weight complexity at that depth, and every
Gaussian exponent $\gamma\ge1+\zeta$. No internal robustness restriction
or infinite-width limit is required. The result covers parity, all
symmetric Boolean functions, and polynomial-size DNF/CNF. An extension
covers targets of any fixed intrinsic depth, with polynomial degree
depending on that intrinsic depth. More precisely, target depth $D_0$
costs a factor $(2+d+C+R)^{O_\zeta(D_0)}$, with polynomial dependence
on every learner depth $L\ge D_0$. The proof combines
a Gaussian-distance decomposition, a width-independent tanh product-norm
bound, and an energy-efficient representation obtained by zero-padding
and scalar transmission. It does not establish the mean-field case
$\gamma=1$ or polynomial dependence on the size of an arbitrary deep
representation.
\end{abstract}

\section{Statement and conventions}

The target $f:\{-1,1\}^d\to\{-1,1\}$ is learned from independent
examples $(x_i,f(x_i))$, where $x_i$ have an arbitrary fixed distribution
$D$. A learner of depth $L\ge2$ has $L-1$ hidden layers, all of width $N$,
and computes
\[
 h^0=(1,x),\qquad h^s=\tanh(W^sh^{s-1})\quad(1\le s<L),
 \qquad z_w(x)=v^Th^{L-1}.
\]
There are no biases after the augmented input. Define
\[
 \kappa_N(w)=\frac dN\|W^1\|_F^2+
    \sum_{s=2}^{L-1}\|W^s\|_F^2+N\|v\|_2^2.
\]
Fix a prescribed scale schedule $t=t(d,L,N)>0$, abbreviated $t(N)$
when $d,L$ are fixed; a constant schedule is allowed.
At each specified dimension, depth and width, the raw Gaussian
$\mu_{N,t}$ has independent entry variances
$t^2/d,t^2/N,t^2/N^2$, respectively. In energy coordinates its density is
proportional to $\exp[-N\kappa_N/(2t^2)]$.
Write $\sigma^2=t^2/N$. The exponent-$\gamma$ convention is
$t^2=N^{1-\gamma}$, equivalently $\sigma^2=N^{-\gamma}$.

The prior conditions this Gaussian on $|z_w(x)|>1/2$ at every Boolean
input. The posterior additionally conditions on
$f(x_i)z_w(x_i)>0$ at every training input. Thus it is equivalently the
raw Gaussian conditioned on the intersection of these events. Throughout,
``unrestricted'' means that no additional internal robustness condition
is imposed; it does not remove the original all-input output margin.

Let $C_L(f)$ be the least width $N$ for which a depth-$L$ representation
has all weights bounded by $N$ and signed output greater than $1/2$
on the entire cube. The main theorem, proved in
Section~\ref{sec:shallow-unrestricted}, states that a width-$C$,
weight-bound-$R$, unit-margin shallow target has sample complexity
\[
 \poly_\zeta(d,C,R,L,\varepsilon^{-1},\log(1/\delta))
\]
for every specified $L\ge2$, simultaneously at all $N\ge C_L(f)$ and
all $0<\sigma^2\le N^{-(1+\zeta)}$. The sample event has probability
at least $1-\delta$, and posterior Gibbs error is at most $\varepsilon$.
The polynomial degree may depend on $\zeta$; no uniform limit as
$\zeta\downarrow0$ is asserted. No claim of computationally efficient
posterior evaluation is made.

For a sample $S$, write $Q_S$ for the weight posterior and define its
Gibbs error by
$\operatorname{GibbsErr}(Q_S)=\mathbb E_{w\sim Q_S}
\Pr_{X\sim D}\{\operatorname{sign}z_w(X)\ne f(X)\}$.
If $p_S(x)$ is posterior probability of the correct label, then
\[
 D\{x:p_S(x)\le q\}
 \le\frac{\operatorname{GibbsErr}(Q_S)}{1-q}\quad(0<q<1).
\]
Hence the same theorem gives the probability-of-correct-answer thresholds
$2/3$ and $1/2$ after multiplying the Gibbs error bound by three and two.
Arbitrary mixtures of tied or untied posterior functions are allowed.
For auxiliary classifiers, take $\operatorname{sign}(0)=+1$.

\paragraph{Proof organization and scope.}
A finite parameter grid handles polynomially bounded learner widths.
At larger widths and $\gamma>1$, a typical Gaussian weight array has small
operator norms in normalized coordinates. Gaussian isoperimetry forces the
posterior close to this typical set. Subtracting the typical part leaves a
low-energy signal with exactly the same Boolean signs. Product-norm
Rademacher control bounds its generalization error. Finally, a shallow
representation can be padded and transmitted through any number of layers
so that the relevant product-norm cost stays polynomial in its shallow
parameters. All arguments and quantitative estimates needed for this
result are included here. These are ordinary mathematical proofs; this
paper does not claim Lean certification.

\paragraph{Related representation-cost questions.}
Jacot~\cite{jacot2023} studies the infinite-depth limit of normalized
representation cost for homogeneous networks. Our scalar padding serves
a related purpose, but tanh is nonhomogeneous and the required estimates
also control finite Gaussian posterior mass at every admissible width.
Depth separations for norm-controlled ReLU classes
\cite{parkinson2024} concern another learning model; they do not by
themselves settle the Gaussian-posterior question considered here.

\section{Unrestricted learning at bounded widths}
\label{sec:finitewidth-unrestricted}

A robustness restriction is unnecessary when the learner's width is itself
polynomially bounded. The following elementary argument retains the original
all-input clean margin and Gaussian conditioning conventions, and is uniform
in the Gaussian scale. It also covers analog target networks at their original
depth, without digital simulation.

Let $P_N=N(d+2)+(L-2)N^2$ be the parameter count, and let $\mu_{N,t}$ denote
the raw Gaussian with variances $t^2/d,t^2/N,t^2/N^2$, respectively. Thus
$\kappa_N=(t^2/N)\chi^2_{P_N}$ under $\mu_{N,t}$. Fix integers
$1\le C\le M$ and suppose, for some $D\ge1$, that
\begin{equation}\label{eq:finitewidth-targetmass}
 \mu_{N,t}\{w:f(x)z_w(x)>1/2\text{ for every }x\}
 \ge \exp(-DN/t^2).
\end{equation}
This holds for $C\le N\le M$ and $0<t\le1$. Only this target-mass
hypothesis is needed; posterior networks may be fragile.

\begin{theorem}\label{thm:finitewidth-unrestricted}
For $0<\varepsilon,\delta<1$, set
\begin{align*}
 B&=4D+2[d+2+(L-2)M]\log2+4\log(2/\varepsilon),\\
 A&=\sqrt{MB},\qquad T=d+1+M(L-1),\qquad
 q=[4T(MA)^{L-1}]^{-1}.
\end{align*}
If
\[
 n\ge\frac2\varepsilon\left[
 \log M+P_M\log(5A/q)+\log(1/\delta)\right],
\]
then, with probability at least $1-\delta$ over independent examples from any
input distribution, the unrestricted Gaussian posterior has Gibbs error at
most $\varepsilon$, simultaneously for every $C\le N\le M$ and $0<t\le1$.
Here ``unrestricted'' retains the original conditioning on
$|z_w(x)|>1/2$ for all inputs, but imposes no internal robustness condition.
The same sample event works after any number of additional correctly labeled
examples. In particular, if $D$ and $M$ are polynomial in $d,L,C$, the sample
bound is polynomial in $d,L,C,\varepsilon^{-1},\log(1/\delta)$, uniformly for
all $\gamma\ge1$ with $t^2=N^{1-\gamma}$.
\end{theorem}

\begin{proof}
For every sample, the event used to condition the raw Gaussian contains the
target event in~\eqref{eq:finitewidth-targetmass}. The chi-square exponential
moment therefore gives
\[
 Q_{S,N,t}\{\kappa_N>B\}
 \le\exp\left(\frac{DN}{t^2}+\frac{P_N}2\log2-
                       \frac{NB}{4t^2}\right)
 \le\varepsilon/2.
\]
The last inequality uses $P_N/N\le d+2+(L-2)M$, $N\ge1$, and $t\le1$.
This bound holds for every sample, without a probabilistic qualification.

If $\kappa_N\le B$, every physical parameter has magnitude at most $A$.
Round each to the nearest multiple of $q$, using at most $5A/q$ grid values
per coordinate. For the maximum coordinatewise hidden-state discrepancies,
boundedness and the Lipschitz property of tanh give
\[
 e_1\le(d+1)q,\qquad e_s\le NAe_{s-1}+Nq,\qquad
 |z-z_{\rm grid}|\le NAe_{L-1}+Nq.
\]
Consequently, uniformly over Boolean inputs,
\[
 |z-z_{\rm grid}|\le q[d+1+M(L-1)](MA)^{L-1}=1/4.
\]
Rounding thus preserves the entire Boolean function of every posterior
network with energy at most $B$.

The union of these grid classifiers over $C\le N\le M$ has logarithmic
cardinality at most $\log M+P_M\log(5A/q)$. A consistency union bound gives
one sample event on which every consistent grid classifier has true error
at most $\varepsilon/2$. Every good-energy posterior draw has the same
function as such a classifier; the remaining posterior mass is at most
$\varepsilon/2$. This proves the Gibbs bound. The finite class is independent
of $t$, and consistency only becomes more restrictive as examples are added.
\end{proof}

The logarithm of the grid size satisfies
\[
 \log(5A/q)=\log(20AT)+(L-1)\log(MA).
\]
The exponentially small parameter precision costs only a polynomial in
$L$ in this logarithm. The unresolved unrestricted problem is therefore
uniformity at learner widths \emph{larger than every polynomial in the
specified target complexity and depth}; this theorem does not show those
larger widths actually require more samples.

\begin{corollary}[Same-depth bounded-weight analog targets]
Suppose a width-$C$, depth-$L$ tanh network realizes $f$ with
$f(x)z(x)\ge1$ on every Boolean input and weights of magnitude at most
$R\ge1$. Fix $M\ge C$ and set
\begin{align*}
 K&=2[d(d+1)+(L-1)C^2]R^2,\\
 a&=[4(d+1+M(L-1))(MR)^{L-1}]^{-1},\\
 D&=\max\{1,K+[d+2+(L-2)M][2+\log(1/a)]\}.
\end{align*}
Then~\eqref{eq:finitewidth-targetmass} holds with this $D$.
Thus Theorem~\ref{thm:finitewidth-unrestricted} supplies a polynomial sample
bound at the original depth for every $C\le N\le M$, with dependence on
$R$ only through $\log R$ in the final sample bound. No robustness condition
on the target implementation or prior support is needed. A target with the
original strict margin $>1/2$ satisfies the displayed hypothesis after
doubling its readout.
\end{corollary}

\begin{proof}
Clone the $C$ neurons into balanced nonempty groups of sizes $n_i$ summing
to $N$. Copy first-layer rows and replace internal and readout weights by
$W_{ij}/n_j$ and $v_i/n_i$. This preserves the exact function. Since
$\max_i n_i/\min_i n_i\le2$, its energy is at most $K$, and its weights
remain bounded by $R$. A physical-coordinate perturbation of magnitude
at most $a$ changes every output by at most $1/4$, by the recurrence above.

Use isotropic coordinates
$y=(\sqrt d\,W^1,\sqrt N\,W^{\rm internal},Nv)$, so
$y\sim\mathcal N(0,t^2I)$ and $\|y_*\|^2\le NK$ for the cloned center.
The standard-normal-coordinate box
$|y_j/t-y_{*,j}/t|\le a$ lies within the physical box, because all coordinate
scalings are at least one and $t\le1$. Its probability is at least
\begin{align*}
 &\left(\frac{2a}{\sqrt{2\pi}}\right)^{P_N}
 \exp\left[-\frac12\sum_j(|y_{*,j}|/t+a)^2\right]\\
 &\qquad\ge\exp[-NK/t^2-P_N(2+\log(1/a))]
 \ge\exp(-DN/t^2).
\end{align*}
Here $(u+a)^2\le2u^2+2a^2$ and $a\le1$ suffice. Every network in this box
has target margin at least $3/4$.
\end{proof}

This is an elementary finite-parameter learning result, rather than a new
claim about the general VC dimension of neural networks. Its role is to
isolate which width regime remains open under the exact prior conventions.

\section{A width-independent tanh bound without robustness}
\label{sec:productnorm}

We first establish the width-independent function-class bound used to
control the low-energy signal extracted from a posterior draw. The proof
accounts for the fact that tanh is not positively homogeneous.

The normalized-activation and exponential-moment method is closely related
to Nakhleh--Nowak's treatment of deep neural variation spaces
\cite{nakhleh2026}, which explicitly uses $\tanh(st)/s$ and slope
discretization, and to the exponential peeling argument of
Golowich--Rakhlin--Shamir~\cite{golowich2018}. We give a self-contained
product-Frobenius version for the present Gaussian-energy conventions;
we do not claim the underlying proof technique is new.

\begin{theorem}[Product-Frobenius bound for tanh]\label{thm:tanh-product}
Let $\mathcal F_{L,P}$ contain all scalar-output tanh networks with $L-1$
hidden layers, arbitrary finite widths, a linear readout, and product of
layer Frobenius norms at most $P$. Biases are absent except through an
augmented input. On every sample $x_1,\ldots,x_n$ with $\|x_i\|\le R$,
\[
 \widehat{\mathfrak R}_n(\mathcal F_{L,P})
 \le \frac{PR}{\sqrt n}
 \left[2+\sqrt{2(L-1)\log\bigl(3+n(L-1)^2\bigr)}\right].
\]
The bound is independent of every hidden width and of the individual layer
norms apart from their product.
\end{theorem}

\begin{proof}
If $R=0$, every sampled input and every network output on the sample is
zero, so the bound is immediate. Assume $R>0$ below.
Define $\phi_s(u)=\tanh(su)/s$ for $s>0$, $\phi_0(u)=u$, and
$\phi_\infty(u)=0$. These functions are odd, $1$-Lipschitz, and zero at zero.
For finite $s,t\ge0$ and $r\ge0$,
\[
 \phi_s(ru)=r\phi_{sr}(u),\qquad
 |\phi_s(u)-\phi_t(u)|\le u^2|s-t|,
 \qquad |\phi_s(u)|\le1/s\quad(s>0).
\]
The derivative bound follows by differentiating
$\phi_s(u)=u\int_0^1\operatorname{sech}^2(suv)\,dv$.

Let $\mathcal G_0$ contain the identity feature map. Inductively allow
arbitrary finite-dimensional features
\[
 g_j(x)=\phi_{s_j}(a_j^T h(x)),\quad
 h\in\mathcal G_{\ell-1},\quad \sum_j\|a_j\|^2\le1,
 \quad s_j\in[0,\infty].
\]
All these feature vectors have norm at most $R$ on the sample. Put
$B_\ell(\epsilon)=\sup_{g\in\mathcal G_\ell}
 \|\sum_i\epsilon_i g(x_i)\|$ for independent Rademacher signs.
Writing $a_j=r_j u_j$ and using the scaling identity shows
\[
 B_\ell\le
 \sup_{h,\,\|u\|\le1,\,s\in[0,\infty]}
 \left|\sum_i\epsilon_i\phi_s(u^T h(x_i))\right|.
\]
For any $\eta>0$, uniformly on $[-R,R]$ the slope family admits an
$\eta/n$ approximation using at most
$J=3+n^2R^2/\eta^2$ choices: truncate $s>n/\eta$ to the zero function,
and round slopes below that threshold down at mesh $\eta/(nR^2)$.
The error in the Rademacher sum is at most $\eta$.

We use the convex contraction principle
\[
 \mathbb E G\!\left(\sup_{t\in T}\sum_i\epsilon_i\psi_i(t_i)\right)
 \le
 \mathbb E G\!\left(\sup_{t\in T}\sum_i\epsilon_i t_i\right)
\]
for convex nondecreasing $G$ and $1$-Lipschitz $\psi_i$ vanishing at zero.
For completeness, the single-coordinate proof compares
$X_\pm=\sup_t(a_t\pm\psi(t_k))$ and
$Y_\pm=\sup_t(a_t\pm t_k)$. One has
$\max(Y_+,Y_-)\ge\max(X_+,X_-)$ by $|\psi(t)|\le|t|$.
An approximately maximizing pair for $X_+,X_-$ also gives
$Y_++Y_-\ge X_++X_-$, by the Lipschitz property.
Convex monotonicity therefore gives
$G(Y_+)+G(Y_-)\ge G(X_+)+G(X_-)$; iterate over coordinates.
Bounded separable classes follow by approximation from finite classes.
In our application, $u$ ranges symmetrically and $\phi_s$ is odd, so the
absolute supremum is the same as the one-sided supremum.

Apply this principle with $G(t)=e^{\lambda t}$ to each grid slope, then
bound the exponential of a maximum by the sum of the exponentials:
\[
 \mathbb E e^{\lambda B_\ell}
 \le e^{\lambda\eta}J\,\mathbb E e^{\lambda B_{\ell-1}}.
\]
Furthermore, $\mathbb E B_0\le R\sqrt n$, and bounded differences gives
$\mathbb E e^{\lambda B_0}
 \le e^{\lambda R\sqrt n+\lambda^2nR^2/2}$.
Writing $H=L-1$, iteration and Jensen yield
\[
 \mathbb E B_H\le H\eta+R\sqrt n+
             \frac{H\log J}{\lambda}+\frac{\lambda nR^2}{2}.
\]
Take $\eta=R\sqrt n/H$ and optimize $\lambda$ to obtain
$\mathbb E B_H\le R\sqrt n[2+\sqrt{2H\log(3+nH^2)}]$.

Finally, for a particular original network, write
$M_\ell=\|W^\ell\|_F$ and $S_\ell=\prod_{j\le\ell}M_j$.
Dividing its $\ell$th hidden vector by $S_\ell$ changes its recursion to
$\phi_{S_\ell}((W^\ell/M_\ell)g^{\ell-1})$.
Thus its normalized hidden features belong to the enlarged classes above,
and its scalar output is at most $P$ times a norm-one linear readout of
$\mathcal G_H$. Zero-norm layers give the zero function and require no
normalization. This proves the theorem. All suprema are measurable by
countable approximation of finite-dimensional parameter spaces and then
union over widths.
\end{proof}

For the Gaussian energy in our model, the product of the $L$ weighted
squared layer norms is $dP^2$. Consequently
\[
 P\le d^{-1/2}(\kappa_N/L)^{L/2}
 \le d^{-1/2}e^{B/2}\qquad\text{if }\kappa_N\le L+B.
\]
Applying the theorem to the ramp loss
$\ell(yz)=\min\{1,\max\{0,1-2yz\}\}$ gives the following explicit bound.

\begin{corollary}[All widths at bounded excess energy]\label{cor:excess-risk}
Fix $B\ge0$. With probability at least $1-\delta$, every sample-consistent
depth-$L$ network, at any finite width, with training margins at least $1/2$
and energy $\kappa_N\le L+B$ has sign error at most
\[
 \frac{4\sqrt2 e^{B/2}}{\sqrt n}
 \left[2+\sqrt{2(L-1)\log\bigl(3+n(L-1)^2\bigr)}\right]
 +\sqrt{\frac{\log(1/\delta)}{2n}}.
\]
No internal robustness assumption is imposed.
\end{corollary}

\begin{proof}
Augmented Boolean inputs have norm $\sqrt{d+1}$, so $PR\le\sqrt2e^{B/2}$.
The ramp loss is bounded by one and is $2$-Lipschitz. Symmetrization,
scalar contraction, and bounded differences give population loss at most
empirical loss plus four times the network Rademacher bound and the
stated confidence term. Empirical loss is zero and population loss
upper-bounds sign error.
\end{proof}

\begin{lemma}[A scalar transmission construction]\label{lem:scalar-transmission}
For every real $D>0$ and integer $m\ge2$, there are positive scalar weights
$\alpha,b_2,\ldots,b_{m-1},c$ such that
\[
 r_1=\tanh\alpha,\quad r_s=\tanh(b_s r_{s-1}),\quad
 cr_{m-1}=1/2,
\]
and, writing $D'=\max\{D,4\}$,
\[
 D\alpha^2+\sum_{s=2}^{m-1}b_s^2+c^2
 \le(m+12)(D'/4)^{1/m},\qquad c^2=(D'/4)^{1/m}.
\]
\end{lemma}

\begin{proof}
Set $q=(D'/4)^{1/(2m)}$, $k_s=\min\{s,m-s\}$, and
$r_s=q^s/(\sqrt{D'}k_s)$. Then $0<r_s\le1/(2qk_s)\le1/2$ and
$r_{m-1}=1/(2q)$. Choose
$\alpha=\operatorname{atanh}r_1$,
$b_s=\operatorname{atanh}(r_s)/r_{s-1}$ and $c=q$.
For $0<u\le1/2$ the power series gives
$(\operatorname{atanh}(u)/u)^2\le1+u^2$.
Writing $R_s=k_{s-1}/k_s$, the reciprocal pairs
$k/(k+1),(k+1)/k$ satisfy
\[
 \left(\frac{k+1}{k}\right)^2+
 \left(\frac{k}{k+1}\right)^2-2\le4/k^2.
\]
Consequently $\sum_{s=2}^{m-1}R_s^2\le m-2+7$,
$R_s\le2$, and $\sum_{s=1}^{m-1}k_s^{-2}<4$.
Using $q^{2m}=D'/4$ therefore gives
\[
 D\alpha^2\le q^2+1/4,\qquad
 \sum_{s=2}^{m-1}b_s^2\le q^2(m-2+7)+4,\qquad c^2=q^2.
\]
Their sum is at most $(m+12)q^2$. Empty internal sums at $m=2$
satisfy the same estimates.
\end{proof}


\section{Unrestricted finite-variance learning for \texorpdfstring{$\gamma>1$}{gamma greater than one}}
\label{sec:gaussian-distance}

There is a stronger bridge than concentration of the full Frobenius energy.
A typical Gaussian bulk has small operator norms when $\gamma>1$, even when
its total Frobenius energy is large. Gaussian isoperimetry shows that the
posterior lies near such a bulk, with the remaining signal having energy
close to a target witness's. Removing the bulk then preserves all Boolean
signs under an explicit product-norm condition. No additional restriction of prior
support is used.

Write $\sigma^2=t^2/N$, so the original Gaussian has independent entry
variances $N\sigma^2/d,\sigma^2,\sigma^2/N$. Its energy coordinates are
\[
 A_1=\sqrt{d/N}\,W^1,\qquad A_s=W^s\ (2\le s<L),
 \qquad A_L=\sqrt N\,v,
\]
and all their entries are independent $\mathcal N(0,\sigma^2)$.
Their squared Euclidean norm is exactly $\kappa_N$.

\begin{theorem}[Unrestricted learning above exponent one]
\label{thm:unrestricted-above-one}
Fix integers $d\ge1$, $L\ge2$, $C\ge1$, and a real $K\ge L$.
Use tanh hidden activations and a linear readout. Suppose that for every
$N\ge C$ there is a width-$N$, depth-$L$ target witness satisfying
\[
 f(x)z(x)\ge1\quad\hbox{for every Boolean input},\qquad \kappa_N\le K.
\]
Put
\[
 Q_* = \left(\frac{K+2}{L}\right)^{L/2}.
\]
For every fixed $\zeta>0$, every $0<\varepsilon,\delta<1$, and every
distribution $D$ on the Boolean cube, the original Gaussian learner, conditioned only
on the original global clean margin and noiseless training consistency,
has a sample bound polynomial in
\[
 d,L,C,K,Q_*,\varepsilon^{-1},\log(1/\delta)
\]
which, with probability at least $1-\delta$ over independent training
inputs drawn from $D$ and labeled by $f$, gives Gibbs error at most
$\varepsilon$, simultaneously for all
$N\ge C$ and all $0<\sigma^2\le N^{-(1+\zeta)}$.
Thus it covers all $\gamma\ge1+\zeta$ at their actual finite Gaussian scales.
The polynomial degree may depend on $\zeta$; no uniform claim as
$\zeta\downarrow0$ is made.
\end{theorem}

We prove the statement with explicit width thresholds and sample inequalities.
The hypothesis $K\ge L$ costs nothing: averaging over the cube and using
AM--GM shows that a global unit-margin witness has
$\kappa_N\ge Ld^{1/L}\ge L$.

\paragraph{A perturbation estimate using operator norms.}
Normalize hidden states by $\sqrt N$ and the augmented input by $\sqrt d$.
The input norm is at most $\sqrt2$, all coordinate activations remain
$1$-Lipschitz and zero-preserving, and the layer maps are the $A_s$ above.
If $U$ has energy at most $K+1$, and an additive weight perturbation $B$
has operator norm at most $\eta$ in every normalized layer, telescoping
one layer at a time gives
\begin{equation}\label{eq:operator-perturbation}
 \sup_x|z_{U+B}(x)-z_U(x)|
 \le\sqrt2\,L\eta
 \left[\frac{K+1+2\eta\sqrt{L(K+1)}+L\eta^2}{L-1}\right]^{(L-1)/2}.
\end{equation}
Indeed, each summand is at most $\sqrt2\eta$ times a product of $L-1$
other layer norms. The sum of squares of all perturbed layer-norm bounds
is at most the numerator, and AM--GM bounds each such product.
Set
\[
 A=\left(\frac{K+2}{L-1}\right)^{(L-1)/2},\qquad
 \eta=\min\left\{1,\frac1{4\sqrt{L(K+1)}},
                         \frac1{8\sqrt2\,LA}\right\}.
\]
The numerator in~\eqref{eq:operator-perturbation} is at most $K+2$, so the
output change is at most $1/8$. Furthermore
$A\le\sqrt e\,Q_*$, hence $\eta^{-1}$ is polynomial in $L,K,Q_*$.

\paragraph{A symmetric typical Gaussian set.}
For completely explicit constants, put
\[
 a_d^2=2(d+1)\log2+4\log8,\qquad
 b_L^2=16\log9+8\log(8L),\qquad b=\max\{a_d,b_L\}.
\]
Let $T$ consist of bulk matrices whose first and last energy-coordinate
Frobenius norms are at most $ta_d$, and whose internal operator norms are
at most $tb_L$. Then $T$ is closed, convex and centrally symmetric, and
\begin{equation}\label{eq:typical-half}
 \mu_{N,t}(T)\ge1/2.
\end{equation}
For the first block, its squared Frobenius norm is
$(t^2/N)\chi^2_{N(d+1)}$; the readout is the same expression with $d+1$
replaced by one. The exponential moment of $\chi^2$ bounds each failure
probability by $1/8$. For an internal matrix $tG/\sqrt N$, a $1/4$-net of
the unit sphere has at most $9^N$ points, and the two-net scalar Gaussian
bound gives
\[
 \Pr\{\|tG/\sqrt N\|_{\rm op}>tb_L\}
 \le2\exp(2N\log9-Nb_L^2/8)\le\frac1{4L}.
\]
Union bounding over the $L-2$ internal blocks proves~\eqref{eq:typical-half}.
Every normalized block of a member of $T$ has operator norm at most $tb$.

\paragraph{Target mass by a Gaussian shift.}
Suppose $tb\le\eta$. For a target witness $w_*$, the shifted set $w_*+T$
has target margin at least $7/8$, by~\eqref{eq:operator-perturbation}.
Central symmetry and the Gaussian change-of-mean identity yield
\begin{align*}
 \mu(w_*+T)
 &=e^{-\kappa_N(w_*)/(2\sigma^2)}
   \int_T e^{-\langle w_*,B\rangle_\kappa/\sigma^2}\,d\mu(B)\\
 &\ge e^{-K/(2\sigma^2)}\mu(T)
 \ge\tfrac12 e^{-K/(2\sigma^2)}.
\end{align*}
The inequality follows by pairing $B$ with $-B$ and using $\cosh u\ge1$.
This lower bound holds for the raw probability of the conditioning event
for every correctly labeled sample.

\paragraph{The posterior is close to the typical set.}
Let $d_\kappa(w,T)$ be Euclidean distance in energy coordinates. Gaussian
isoperimetry~\cite{borell1975} and~\eqref{eq:typical-half} imply
\[
 \mu\{d_\kappa(w,T)>r\}\le e^{-r^2/(2\sigma^2)}.
\]
Dividing by the preceding target-mass lower bound gives, for every sample,
\begin{equation}\label{eq:posterior-distance-tail}
 Q_s\{d_\kappa(w,T)^2>K+1\}\le2e^{-1/(2\sigma^2)}.
\end{equation}
On the complement, let $B$ be the metric projection of $w$ onto $T$ and
write $U=w-B$. The projection is well defined and continuous, since $T$
is nonempty, closed and convex. Then $\kappa_N(U)\le K+1$ and
$\sup_x|z_w(x)-z_U(x)|\le1/8$. The original posterior has clean margin
$>1/2$ everywhere, so $U$ has \emph{exactly the same Boolean function}
as $w$ and training margins $>3/8$. It is enough to use a margin-$1/4$
ramp loss below.

\paragraph{A polynomial width split.}
Take $\beta=\varepsilon/2$ and
\[
 M=\left\lceil\max\left\{C,2,(b/\eta)^{2/\zeta},
                                      2\log(2/\beta)\right\}\right\rceil.
\]
For every $N\ge M$ and $\sigma^2\le N^{-(1+\zeta)}$, one has
$t\le N^{-\zeta/2}$, hence $tb\le\eta$, and
$2e^{-1/(2\sigma^2)}\le\beta$.
The product norm of $U$ is at most
$d^{-1/2}[(K+1)/L]^{L/2}$. Write
\[
 G=\sqrt{(d+1)/d}\left(\frac{K+1}{L}\right)^{L/2},\qquad H=L-1.
\]
The product-norm Rademacher bound, now with the $4$-Lipschitz ramp loss
$\min\{1,\max\{0,1-4yz\}\}$, gives one sample event of probability at
least $1-\delta/2$ on which every such consistent residual network has
error at most
\begin{equation}\label{eq:distance-sample-bound}
 F_n=\frac{8G}{\sqrt n}\left[2+\sqrt{2H\log(3+nH^2)}\right]
                     +\sqrt{\frac{\log(2/\delta)}{2n}}.
\end{equation}
This event already includes all finite widths. When $F_n\le\varepsilon/2$,
adding the posterior exceptional mass $\beta$ proves the required bound
for every $N\ge M$ and every stated scale.

For $C\le N\le M$, use Theorem~\ref{thm:finitewidth-unrestricted}.
Here is a target-mass constant requiring only the present witness assumption.
Every witness weight is at most $R=\sqrt{MK}$. Put
\begin{align*}
 a&=[4(d+1+M(L-1))(MR)^{L-1}]^{-1},\\
 D_M&=\max\{1,K+[d+2+(L-2)M][2+\log(1/a)]\}.
\end{align*}
The same Gaussian-coordinate box calculation as in the finite-width proof
has physical perturbation at most $a$, retains target margin at least $3/4$,
and has mass at least $\exp(-ND_M/t^2)$.
Applying that theorem with error $\varepsilon$ and failure probability
$\delta/2$ controls all these widths and all $t\le1$ simultaneously.
Its finite-grid logarithm is polynomial in $M,d,L,\log K$.
The union of the two sample events proves the theorem.
Finally, $G\le\sqrt2 Q_*$, $b^2=O(d+\log L)$ and
$A\le\sqrt e Q_*$ show that $M$ and a sample size satisfying all the
displayed inequalities are polynomial in the stated arguments for fixed
$\zeta$. This proves the asserted quantitative dependence.

\paragraph{Scope.}
This argument does not apply at $\gamma=1$: the Gaussian bulk's operator
norm does not tend to zero with width there. It also does not make $Q_*$
polynomial for every small circuit. If a fixed base witness can be padded
with energy $K=L+B$, where $B$ is independent of the added depth, then
$Q_*\le e^{(B+2)/2}$. The theorem therefore removes exponential dependence
on \emph{additional depth}, at every width and every fixed exponent above
one, while retaining a potentially exponential dependence on base-witness
energy. Stronger target-specific energy estimates can make $Q_*$ itself
polynomial.

\section{Polynomial unrestricted learning of shallow targets at every depth}
\label{sec:shallow-unrestricted}

A shallow target can have a small first-layer Gaussian energy at a
polynomial base width simply by adding zero neurons. Its readout energy
is then polynomial. Distributing that readout cost over the learner's
remaining layers makes the product-norm parameter polynomial uniformly
in depth. This gives a family for which neither arbitrary width nor
arbitrary added depth requires a robustness restriction.

\begin{theorem}[Shallow targets, all depths and all admissible widths]
\label{thm:shallow-all-depths}
Suppose a Boolean target $f:\{-1,1\}^d\to\{-1,1\}$ has a one-hidden-layer
representation
\[
 z(x)=\sum_{j=1}^{C}v_j\tanh(a_j^T\bar x),\qquad
 \bar x=(1,x),\qquad f(x)z(x)\ge1
\]
with every coordinate of $a_j$ and every $v_j$ bounded in absolute value
by $R\ge1$. Let $C_L(f)$ denote its bounded-weight network complexity
at depth $L$: the least width $N$ admitting a realization with all
weights at most $N$ in absolute value and signed output greater than
$1/2$ everywhere.

For each fixed $\zeta>0$, the original unrestricted Gaussian learner has
sample complexity polynomial in
\[
 d,C,R,L,\varepsilon^{-1},\log(1/\delta),
\]
for every specified learner depth $L\ge2$ and every input distribution,
simultaneously for every width $N\ge C_L(f)$ and every
$\gamma\ge1+\zeta$. With probability at least
$1-\delta$ over a sample of that size, posterior Gibbs error is at most
$\varepsilon$ simultaneously at all these widths and all prior scales
$0<\sigma^2\le N^{-(1+\zeta)}$. The polynomial's degree may depend on
$\zeta$. The prior is conditioned only on the original global clean
margin and the training labels.
\end{theorem}

\begin{proof}
Let $A$ have rows $a_j^T$, and set
\[
 M=\left\lceil\max\{4,C,2R,d\|A\|_F^2\}\right\rceil,
 \qquad T=\max\{1,M\|v\|_2^2\}.
\]
Both are polynomial in $d,C,R$. Add zero hidden neurons to obtain width
$M$. This preserves the function, and its first-layer and readout
energies satisfy
\[
 E_{\rm pre}=\frac dM\|A\|_F^2\le1,
 \qquad c^2=M\|v\|_2^2\le T.
\]
In fact $M\|v\|_2^2\ge1$, since unit margin and $|\tanh|\le1$ imply
$\|v\|_2\ge1/\sqrt C$, but the maximum makes this unneeded.

For any $N\ge M$, balanced cloning of this shallow network preserves its
output exactly and gives
\[
 E_{\rm pre}\le2,\qquad c^2=N\|v_N\|_2^2\le2T.
\]
Explicitly, assign each of the $M$ original neurons either
$\lfloor N/M\rfloor$ or $\lceil N/M\rceil$ clones, copy its incoming
weights, and divide its readout weight by its number of clones.
The first cost increases by at most $M\lceil N/M\rceil/N\le2$,
and the readout cost by at most $N/(M\lfloor N/M\rfloor)\le2$.
For learner depth $L=2$ we already have a unit-margin witness with
\[
 K_2=2+2T,\qquad
 Q_*^2=\left(\frac{K_2+2}{2}\right)^2=(T+2)^2.
\]

For $L\ge3$, apply Lemma~\ref{lem:scalar-transmission} with
$D=4c^2$ and $m=L-1$. Replace the old readout by a layer whose every
row is $2\alpha v_N^T$. Follow it with the scalar transmission layers,
implemented by equal coordinates with dense matrix entries $b_s/N$,
and final readout entries $q/N$, where
\[
 q^2=\max\{c^2,1\}^{1/(L-1)}\le(2T)^{1/(L-1)}.
\]
The target had margin at least one, hence in particular at least
$1/2$, so the scalar construction gives margin at least $1/2$.
Double its final readout to restore margin at least one. This adds
$3q^2$ to the energy. Thus every width $N\ge M$ has a depth-$L$
unit-margin witness with
\[
 \kappa_N\le K_L:=2+(L+14)(2T)^{1/(L-1)}.
\]
Writing $q_*^2=(2T)^{1/(L-1)}\ge1$, we have
$K_L+2\le(L+18)q_*^2$. Consequently
\begin{equation}\label{eq:shallow-product-polynomial}
 Q_*^2=\left(\frac{K_L+2}{L}\right)^L
 \le e^{18}(2T)^{L/(L-1)}
 \le e^{18}(2T)^{3/2}\qquad(L\ge3).
\end{equation}
The exponent $L/(L-1)$ is at most $3/2$ for $L\ge3$.
Theorem~\ref{thm:unrestricted-above-one}, with base width $M$, therefore
provides a sample bound polynomial in the stated quantities for all
$N\ge M$.

We verify admissible smaller widths rather than silently identifying
$M$ with the target's minimum complexity. First, $C_L(f)\le M$ for all
$L\ge2$. At width $M$, the original incoming weights have magnitude
at most $R\le M/2$. In the added scalar channel,
$\alpha\le\operatorname{atanh}(1/2)<0.55$, so its first weights
$2\alpha v_j$ have magnitude at most $1.1R<M$.
The scalar construction has $b_s\le2.2q$ and final weights $2q/M$.
Since $T\le MC R^2\le M^4$ and $L\ge3$,
\[
 q\le(2T)^{1/4}\le2^{1/4}M,
 \qquad b_s/M<3\le M,\qquad 2q/M<3\le M.
\]
Thus this depth-$L$, width-$M$ circuit satisfies the bounded-weight
convention and has margin at least one. The case $L=2$ is immediate.

For $C_L(f)\le N<M$, start from a defining width-$C_L(f)$ circuit,
clone it to width $N$, and double its readout. It has margin greater
than one, all weights at most $2M$, and Gaussian energy bounded by a
polynomial in $d,L,M$. The Gaussian box estimate and
Theorem~\ref{thm:finitewidth-unrestricted} consequently give a sample
bound polynomial in $d,L,M,\varepsilon^{-1},\log(1/\delta)$,
simultaneously for all these widths and all $t\le1$.
Use failure probability $\delta/2$ for this event and $\delta/2$ for
the preceding event for $N\ge M$. Their intersection proves the claim.
\end{proof}

The complexity parameter controlling this theorem's sample size is the
size and weight bound of a \emph{shallow} representation. Its width
admissibility is the actual depth-$L$ complexity $C_L(f)$. The theorem
does not identify shallow complexity with the minimum complexity of an
arbitrary deep circuit, and therefore does not settle the unrestricted
problem for all small deep circuits. It also makes no claim at $\gamma=1$.

\begin{corollary}[Parity at arbitrary depth and admissible width]
\label{cor:parity-full-polynomial}
For every fixed $\zeta>0$, parity of any $k\ge1$ chosen coordinates in
$\{-1,1\}^d$ has unrestricted Gaussian posterior sample complexity
polynomial in $d,k,L,\varepsilon^{-1},\log(1/\delta)$, for every
$L\ge2$, every $N\ge C_L(f)$, and every $\gamma\ge1+\zeta$.
The bound holds for every input distribution and for the original
finite-variance posterior.
\end{corollary}

\begin{proof}
Let $q=(k-\sum_{i\in S}x_i)/2\in\{0,\ldots,k\}$ count negative
coordinates, and put $a=\log(16k)$. Define
\[
 \tau_j(x)=\tanh\bigl(a(q-j+1/2)\bigr),\qquad 1\le j\le k.
\]
Every argument has magnitude at least $a/2$, so
$|\tau_j-\operatorname{sign}(q-j+1/2)|\le2e^{-a}=1/(8k)$.
Writing $c_k=1$ for even $k$ and $c_k=0$ for odd $k$, the exact step
identity is
\[
 (-1)^q=c_k+\sum_{j=1}^k(-1)^j
                  \operatorname{sign}(q-j+1/2).
\]
Hence $F=c_k+\sum_j(-1)^j\tau_j$ satisfies
$|F-(-1)^q|\le1/8$. Represent the constant by one extra hidden neuron
with augmented-input weight $\operatorname{atanh}(1/2)$ and no other
incoming weights. A linear readout equal to $2F$ has signed margin
at least $7/4$. Its threshold-neuron readout weights are $2(-1)^j$,
and its constant-neuron weight is $4c_k$.
Its incoming weights have magnitude at most $a(k+1)$.
Thus $C=k+1$ and $R=\max\{4,a(k+1)\}$ are polynomial in $k$.
Apply Theorem~\ref{thm:shallow-all-depths}.
\end{proof}

\begin{corollary}[Symmetric Boolean targets]
The conclusion of Corollary~\ref{cor:parity-full-polynomial} holds for
any Boolean function of the Hamming weight of $k$ selected coordinates.
\end{corollary}

\begin{proof}
If its values are $h(0),\ldots,h(k)\in\{-1,1\}$, then
\[
 h(q)=\frac{h(0)+h(k)}2+
       \sum_{j=1}^k\frac{h(j)-h(j-1)}2
                    \operatorname{sign}(q-j+1/2).
\]
Each coefficient has magnitude at most one. Replacing the signs by
the same $\tau_j$ and doubling the output gives margin at least $7/4$,
using the same width and weight bounds as in the parity proof.
\end{proof}


\begin{corollary}[Polynomial-size DNF and CNF targets]
A Boolean function on $d$ bits represented by a DNF or CNF with $s\ge1$
clauses has unrestricted Gaussian posterior sample complexity polynomial
in $d,s,L,\varepsilon^{-1},\log(1/\delta)$, for every specified $L\ge2$,
every $N\ge C_L(f)$, and every $\gamma\ge1+\zeta$, with $\zeta>0$ fixed.
\end{corollary}

\begin{proof}
For a DNF, let $q_j(x)$ be the number of false literals in clause $j$,
and put $a=\log(16s)$ and
$\tau_j(x)=\tanh(a(1/2-q_j(x)))$.
The exact clause indicator is
$[1+\operatorname{sign}(1/2-q_j)]/2$.
Thus
\[
 H(x)=s-1+\sum_{j=1}^s\operatorname{sign}(1/2-q_j(x))
\]
is $-1$ when no clause is satisfied and at least $1$ otherwise.
Replacing the signs by $\tau_j$ changes $H$ by at most $1/8$.
Doubling the resulting expression gives global signed margin at least
$7/4$. It is a shallow tanh network with $s+1$ features: the extra
feature is the constant $1/2$, with readout weight $4(s-1)$.
After removing duplicate literals and discarding contradictory clauses,
each clause has at most $d$ literals. Each $q_j$ is affine in the bits,
so the remaining incoming and readout weights have magnitude at most
$R=\max\{4s,a(d+1)\}$. The shallow-target theorem applies.
A CNF is the negation of a DNF in negated literals, and negating the
readout preserves every bound. Constant functions are immediate.
\end{proof}

\subsection{Targets of any fixed intrinsic depth}
\label{sec:fixed-intrinsic-depth}

The shallow-target argument extends to any fixed target depth by using
different numbers of clones at its successive hidden layers. The degree
of the resulting polynomial may depend on the target depth, but not on
the learner's additional depth.

\begin{corollary}[Fixed target depth, arbitrary learner depth]
\label{cor:fixed-intrinsic-depth}
Fix an integer $D_0\ge2$ and $\zeta>0$. Suppose $f$ has a width-$C$,
depth-$D_0$ tanh realization in the present architecture, with signed
output at least one on the Boolean cube and all weights bounded by
$R\ge1$. For every specified learner depth $L\ge D_0$, the original
unrestricted Gaussian learner has sample complexity polynomial in
\[
 d,C,R,L,\varepsilon^{-1},\log(1/\delta).
\]
The degree and constants may depend on $D_0$ and $\zeta$.
For every input distribution, one sample event of probability at least
$1-\delta$ gives Gibbs error at most $\varepsilon$ simultaneously for
every $N\ge C_L(f)$ and every
$0<\sigma^2\le N^{-(1+\zeta)}$.
Here $C_L(f)$ is the actual bounded-weight network complexity defined in
Theorem~\ref{thm:shallow-all-depths}. The prior retains the original
global clean-margin conditioning and sign-only training likelihood.
\end{corollary}

\begin{proof}
Put $H=D_0-1$. Pad the target's hidden layers to width $C$ with zero
neurons if necessary, and write $W^1,\ldots,W^H,v$ for its weights.
Choose an integer
\[
 B=\left\lceil\max\left\{4,2R,\frac dC\|W^1\|_F^2,
                 \|W^2\|_F^2,\ldots,\|W^H\|_F^2\right\}\right\rceil,
 \qquad M=CB^H,
\]
omitting the internal-layer terms when $H=1$.
Since $\|W^1\|_F^2\le C(d+1)R^2$ and
$\|W^s\|_F^2\le C^2R^2$ internally, $B$ is polynomial in $d,C,R$,
and so is $M$ for fixed $D_0$.

At hidden layer $s$, give each original neuron exactly
\[
 n_s=B^{H-s}
\]
active copies, and fill the other positions up to width $M$ with zero
neurons. Copy first-layer rows. For $s\ge2$, give every edge from a copy
of source $j$ to a copy of target $i$ weight $W^s_{ij}/n_{s-1}$.
All unused edges are zero. Since $n_H=1$, retain the original readout
on the final active neurons. This preserves the function exactly.
The first-layer and internal energy terms are respectively
\[
 \frac dM n_1\|W^1\|_F^2
      =\frac d{CB}\|W^1\|_F^2\le1,
 \qquad
 \frac{n_s}{n_{s-1}}\|W^s\|_F^2
      =\frac1B\|W^s\|_F^2\le1.
\]
Thus the energy before the readout is at most $H$. Put
\[
 T=\max\{1,M\|v\|_2^2\}\le MC R^2.
\]
It is polynomial in $d,C,R$ for fixed $D_0$.

Balanced cloning of this width-$M$ network into any $N\ge M$ preserves
the function and gives
\[
 E_{\rm pre}\le2H,\qquad c^2=N\|v_N\|_2^2\le2T.
\]
If $L=D_0$, these are already unit-margin witnesses with
$K=2H+2T$, and
\[
 Q_*^2=\left(\frac{K+2}{D_0}\right)^{D_0}
       \le(4T)^{D_0}.
\]

If $L>D_0$, set $m=L-H\ge2$. Replace the old readout by the scalar
transmission from Lemma~\ref{lem:scalar-transmission}, using $D=4c^2$,
and double the final scalar readout to obtain unit margin. Its readout
cost before doubling is
$q^2=\max\{c^2,1\}^{1/m}\le(2T)^{1/m}$.
The total energy is at most
\[
 K_L=2H+(m+15)(2T)^{1/m}.
\]
This is the desired depth: the $H$ old hidden layers and $m-1$ new hidden
layers total $L-1$. Since $(2T)^{1/m}\ge1$,
\begin{align*}
 Q_*^2
 &=\left(\frac{K_L+2}{L}\right)^L\\
 &\le\left(1+\frac{H+17}{L}\right)^L(2T)^{L/(L-H)}\\
 &\le e^{H+17}(2T)^{(D_0+1)/2}.
\end{align*}
The last step uses $L/(L-H)=1+H/m\le1+H/2$.
For $L\ge2D_0$, the sharper upper bound
$Q_*^2\le e^{H+17}(2T)^2$ also holds.
Both $K_L$ and $Q_*$ are polynomial in the stated parameters when
$D_0$ is fixed, uniformly over the learner depth.
Theorem~\ref{thm:unrestricted-above-one} now controls every $N\ge M$.

It remains to check admissible widths below $M$.
First $C_L(f)\le M$ at every $L\ge D_0$.
The tapered construction has all original weights at most $R\le M/2$.
In a padded scalar channel, $\alpha\le\operatorname{atanh}(1/2)<0.55$,
so its first added weights $2\alpha v_j$ have magnitude less than $M$.
The scalar construction has $b_s<2.2q$, and
\[
 T\le MC R^2\le M^4,\qquad
 q\le(2T)^{1/4}\le2^{1/4}M.
\]
Its dense internal entries $b_s/M$ and final readout entries $2q/M$
are therefore less than $3\le M$.
This gives a bounded-weight width-$M$ witness with unit margin at each
learner depth; for $L=D_0$ no scalar channel is needed.

For $C_L(f)\le N<M$, clone a defining width-$C_L(f)$ realization to
width $N$ and double its readout. Its margin is greater than one and
all weights are at most $2M$. The Gaussian box estimate and
Theorem~\ref{thm:finitewidth-unrestricted} give a bound polynomial in
$d,L,M,\varepsilon^{-1},\log(1/\delta)$ for all these widths and all
$t\le1$. Allocate failure probability $\delta/2$ to this event and to
the preceding event for $N\ge M$. Their intersection proves the claim.
\end{proof}

Thus polynomial dependence on learner depth holds for targets of any
fixed intrinsic depth with polynomial target-size dependence. The width
$M=CB^{D_0-1}$ and the exponents above explain why this does not give a
single polynomial jointly in the intrinsic target depth and target size.
The mean-field exponent $\gamma=1$ is also still outside this result.

\begin{corollary}[Quantitative dependence on intrinsic depth]
\label{cor:intrinsic-depth-quantitative}
Fix $\zeta>0$. There are constants $a_\zeta,b_\zeta>0$, independent of
the intrinsic depth, with the following property. Suppose $f$ has a
width-$C$, depth-$D_0$ tanh realization with unit signed margin and
weights bounded by $R\ge1$, where $D_0\ge2$, and set
$S=2+d+C+R$. For every specified learner depth $L\ge D_0$, the sample
bound
\[
 n\ge\left\lceil S^{a_\zeta D_0}
       \bigl(L+\varepsilon^{-1}+\log(2/\delta)\bigr)^{b_\zeta}
       \right\rceil
\]
suffices for posterior Gibbs error at most $\varepsilon$, with probability
at least $1-\delta$ for every input distribution, simultaneously for all
$N\ge C_L(f)$ and $0<\sigma^2\le N^{-(1+\zeta)}$.
The prior and likelihood are those of
Corollary~\ref{cor:fixed-intrinsic-depth}.

Thus the dependence on intrinsic depth is
$\exp[O_\zeta(D_0\log S)]$, while the dependence on learner depth remains
polynomial with degree independent of $D_0$. In particular, if
$D_0\le c\log S$ for a fixed constant $c$, the target-parameter factor
is at most $\exp[O_{\zeta,c}((\log S)^2)]$.
This is a quasipolynomial bound in $S$, not a single polynomial jointly
in intrinsic depth and target size.
\end{corollary}

\begin{proof}
First suppose $L\ge2D_0$. Use the tapered construction in the preceding proof, with
$H=D_0-1$, $M=CB^H$, and $T=\max\{1,M\|v\|_2^2\}$.
Its explicit definition gives $B\le S^5$, so
$M\le S^{5D_0}$ and $T\le MC R^2\le S^{7D_0}$.
Put $m=L-H$. Since $L\ge2D_0$, one has $L/m\le2$ and
\[
 K_L=2H+(m+15)(2T)^{1/m}=O(LT),\qquad
 Q_*^2\le e^{H+17}(2T)^{L/m}
          \le e^{H+17}(2T)^2=S^{O(D_0)}.
\]

For the remaining range $D_0\le L<2D_0$, do not taper the target.
Instead pad its original realization to
$M_0=\lceil\max\{4,C,2R\}\rceil\le S^2$, and put
$T_0=\max\{1,M_0\|v\|_2^2\}\le S^5$.
Balanced cloning to every $N\ge M_0$ gives
\[
 E_{\rm pre}\le2H S^4,\qquad c^2\le2T_0.
\]
At $L=D_0$, retain the readout. For $L>D_0$, append the same scalar
transmission with $m=L-H\ge2$. These constructions give, respectively,
\[
 K_L\le2H S^4+2T_0,
 \qquad
 K_L\le2H S^4+(m+15)(2T_0)^{1/m}.
\]
In either case $K_L\le19LS^5$. Therefore
\[
 Q_*^2\le(20S^5)^L\le S^{7L}\le S^{14D_0}.
\]
The width-$M_0$ witness has weights bounded by $M_0$: the original
weights satisfy $R\le M_0/2$, and the scalar-channel check from the
preceding proof applies because $T_0\le M_0^4$.
Thus $C_L(f)\le M_0$ in this range.

In both depth ranges, the polynomial in
Theorem~\ref{thm:unrestricted-above-one} has degree and constants
depending only on $\zeta$. Substituting the corresponding bounds gives
the stated sample form above the respective base width $M$ or $M_0$.
For smaller admissible widths, the defining-witness and Gaussian-box
argument in the preceding proof gives a sample bound polynomial in
$d,L,M,\varepsilon^{-1},\log(1/\delta)$ (or with $M_0$ in place of $M$),
with degree independent of $D_0$. Substitution gives the same sample form.
Split the failure probability between the two width ranges and enlarge
$a_\zeta,b_\zeta$ to finish.
\end{proof}


\section{Limitations}

The result separates learner depth from shallow target complexity. A target
that has only an exponentially large shallow representation need not receive
a polynomial bound in the size of a much smaller deep representation.
This paper does not settle the unrestricted problem for arbitrary small
deep circuits. As of 29 October 2026, a later result~\cite{conditional2026}
rules out a universal polynomial bound in target resources and learner
depth, uniformly over admissible widths, assuming nonuniform
$\mathsf P/\mathrm{poly}\ne\mathsf{TC}^1$.
The unconditional classification remains open.
The mean-field exponent $\gamma=1$ also remains outside the Gaussian-distance
argument: its typical hidden Gaussian bulk has order-one operator norm.
This failure of one proof is not a sample lower bound. The present result
retains the original all-input output-margin conditioning; removing that
condition would require another argument or a change of prediction and
training conventions.

\begin{thebibliography}{9}
\bibitem{conditional2026} ChatGPT-6 Astra, K. H\"anni, and D. Vaintrob.
A conditional obstruction to polynomial sample complexity in unrestricted NNbayes.
29 October 2026.
\href{https://github.com/kaarelh/AI-works/tree/main/neural-network-generalization/works-in-progress/conditional-nonpolynomial-learning}{Source and PDF}.
\bibitem{jacot2023} A. Jacot.
Implicit bias of large depth networks: a notion of rank for nonlinear functions.
\emph{ICLR}, 2023. \url{https://arxiv.org/abs/2209.15055}.
\bibitem{parkinson2024} S. Parkinson, G. Ongie, R. Willett, O. Shamir,
and N. Srebro. Depth separation in norm-bounded infinite-width neural networks.
\emph{COLT}, 2024. \url{https://proceedings.mlr.press/v247/parkinson24a.html}.
\bibitem{borell1975} C. Borell.
The Brunn--Minkowski inequality in Gauss space.
\emph{Inventiones Mathematicae} 30, 207--216, 1975.
\url{https://doi.org/10.1007/BF01425510}.
\bibitem{golowich2018} N. Golowich, A. Rakhlin, and O. Shamir.
Size-independent sample complexity of neural networks.
\emph{COLT}, 2018. \url{https://proceedings.mlr.press/v75/golowich18a.html}.
\bibitem{nakhleh2026} J. Nakhleh and R. D. Nowak.
Deep Neural Variation Spaces: A Unifying Perspective on Depth and Complexity.
2026. \url{https://arxiv.org/abs/2607.05546}.
\end{thebibliography}
\end{document}
